\documentclass[11pt]{article}

\usepackage[review]{acl}
\usepackage{times}
\usepackage{latexsym}
\usepackage{booktabs}
\usepackage{multirow}
\usepackage{graphicx}
\usepackage{url}
\usepackage{xcolor}

\newcommand{\dn}{\textsc{DN}}

\title{\dn: A Pre-Generation Based Interactive Narrative System with Full-Ready Multimodal Delivery}

\author{Anonymous EMNLP Submission}

\begin{document}
\maketitle

\begin{abstract}
Interactive narrative systems increasingly combine language generation, state tracking, and visual scene synthesis, but this richer output scope can make player-facing latency a central bottleneck. This draft reports the efficiency component of \dn, a pre-generation based interactive narrative system that targets full-ready multimodal delivery. On a 20-item formal benchmark, simulated 60-second player reading time reduces \dn's mean next-turn full-ready latency from 27.095 seconds to 5.299 seconds and its median latency from 24.571 seconds to 0.027 seconds, while preserving a strict success rate of 1.000. Visual quality, narrative quality, and ablation results will be integrated in subsequent revisions.
\end{abstract}

\section{Introduction}
Interactive narrative systems must balance rich multimodal generation with low user-perceived latency. Unlike single-turn text or image generators, a playable narrative system must maintain story state, expose actionable options, advance the plot after a user decision, and deliver the next piece of content in a form that the player can immediately continue from. When the next turn contains both text and a scene image, the relevant efficiency question is therefore not merely how quickly a model can produce an isolated image, but how quickly the system can make the complete next turn usable.

\dn addresses this problem through a pre-generation based delivery mechanism. Instead of waiting until the user clicks the next action to begin all downstream work, \dn uses the player's reading interval to prepare likely next-turn content in the background. This design turns reading time into computation time and aims to reduce the latency that remains visible to the player. The present draft first stabilizes the efficiency evidence for this mechanism. Later revisions will extend the evaluation to visual quality, narrative/text quality, and component-level ablations, so that the final paper can jointly answer whether \dn is fast enough for interaction and whether its delivered content is good enough for narrative play.

The current efficiency results support three stage-level findings. First, strict full-ready delivery without dwell time remains expensive because the system must wait for both textual and visual readiness. Second, a realistic reading interval substantially changes the latency distribution: in the pregen60 setting, most samples are delivered nearly immediately after the next click. Third, comparisons against image-continuation references should be interpreted as scope-aware references rather than fully isomorphic system comparisons, because \dn returns text plus image plus playable state while the baselines report image or workflow latency.

\section{System Overview}
\dn is an interactive narrative generation system that combines structured story-state management, option generation, and scene-image delivery. At each turn, the system presents narrative text and choices to the player, receives the selected action, updates the story state, and prepares the next multimodal scene. The paper's central efficiency hypothesis is that pre-generation can reduce user-visible delay by using the idle interval between content delivery and the player's next action.

The pre-generation mechanism maintains candidate next-turn work in the background. During the player's reading interval, \dn prepares likely continuations and their associated scene images. When the player later selects an action, the system attempts to serve a strict full-ready result from the prepared content; if the requested content is not yet fully available, the measured path waits until both text and image are ready. This distinction is important for the evaluation: a response is not counted as successful merely because a placeholder or partial output exists.

\section{Experimental Setup}
We evaluate efficiency on a 20-item formal benchmark subset derived from the DN-quality benchmark. Each item defines a narrative theme and constraints that can be mapped either to \dn's full interactive pipeline or to image-generation baseline adapters. The evaluation uses the latest v15 composite20 artifacts as the source of truth and excludes earlier v14, fixed-v2, and trimmed-sample exploratory results from the main paper tables.

The primary comparison is not intended to claim that all systems are architecturally identical. Instead, it measures the latency of \dn's full-ready delivery and reports image-baseline latencies as reference points for the image-continuation subtask. This scope-aware framing is necessary because \dn performs story-state update and textual continuation in addition to scene-image delivery, whereas StoryDiffusion, SDM-v2, and IC-LoRA are evaluated as next-turn image or image-workflow references.

\section{Efficiency Evaluation}
\label{sec:efficiency}

\paragraph{Evaluation goal.}
The efficiency evaluation asks how long a player must wait after issuing a next-turn action before the system can deliver content that is ready to continue. This question is central to playable narrative generation: even high-quality narrative content can break interaction if the user must wait for a long or unpredictable generation process. We therefore evaluate \dn under a \emph{full-ready} latency definition, where a response is considered ready only when both the textual continuation and the corresponding scene image are available.

\paragraph{Protocol.}
We evaluate \dn under two settings. The first setting, \textit{normal20}, is a zero-dwell stress setting: the second click is issued immediately, so the system receives little opportunity to benefit from player reading time. The second setting, \textit{pregen60}, simulates a player spending 60 seconds reading the current turn before clicking the next action. This setting tests the intended use case of pre-generation, where the system uses the player's reading interval to prepare likely next-turn content in the background.

For external references, we compare against StoryDiffusion~\citep{zhou2024storydiffusion}, SDM-v2, which is based on latent diffusion modeling~\citep{rombach2022latentdiffusion}, and IC-LoRA~\citep{huang2024iclora}. These baselines are image-continuation references rather than complete interactive narrative systems: they generate a next-turn image or image workflow output from DN-style prompts, but they do not generate the full textual continuation, update the game state, or return a complete playable state. Consequently, the comparison is deliberately conservative in the sense that \dn is measured on a heavier output scope. We keep this distinction explicit throughout the analysis.

\paragraph{Metrics.}
We report success rate, mean latency, median latency, and 95th-percentile latency (P95). For \dn, latency is measured as \texttt{full\_ready\_elapsed\_s}. For the image baselines, latency is measured as next-turn image or workflow latency. We also report mechanism-specific signals for \dn, including direct image rate, likely hit rate, and placeholder rate. These signals help determine whether pre-generation improves latency by making the next-turn scene image available before the user clicks.

\begin{table*}[t]
\centering
\small
\resizebox{\textwidth}{!}{%
\begin{tabular}{llrrrrr}
\toprule
System & Output Scope & $N$ & Mean (s) & Median (s) & P95 (s) & Success \\
\midrule
DN normal20 & Text + image full-ready & 20 & 27.095 & 24.571 & 46.196 & 1.000 \\
DN pregen60 & Text + image full-ready & 20 & 5.299 & 0.027 & 27.150 & 1.000 \\
StoryDiffusion & Continuation image & 20 & 8.495 & -- & 8.900 & 1.000 \\
SDM-v2 & Continuation image & 20 & 2.827 & -- & 2.831 & 1.000 \\
IC-LoRA & Continuation workflow & 20 & 30.057 & -- & 30.059 & 1.000 \\
\bottomrule
\end{tabular}
}
\caption{Efficiency comparison on the formal20 benchmark. \dn reports full-ready latency, where both the textual continuation and the corresponding scene image are available. Image baselines report next-turn image or workflow latency only.}
\label{tab:efficiency-main}
\end{table*}

\begin{table}[t]
\centering
\small
\resizebox{\columnwidth}{!}{%
\begin{tabular}{lrrrr}
\toprule
DN Setting & Mean & Median & P95 & Direct Image Rate \\
\midrule
normal20 & 27.095 & 24.571 & 46.196 & 1.000 \\
pregen60 & 5.299 & 0.027 & 27.150 & 1.000 \\
\bottomrule
\end{tabular}
}
\caption{Effect of simulated reading time on \dn full-ready latency. The pregen60 setting simulates a 60-second dwell interval before the next click.}
\label{tab:dn-pregen}
\end{table}

\paragraph{Main results.}
Table~\ref{tab:efficiency-main} shows that \dn's zero-dwell full-ready latency is 27.095 seconds on average, with a P95 latency of 46.196 seconds. This setting is intentionally strict: it measures the cost of producing a full next-turn response when the system is given essentially no reading interval in which to complete pre-generation. Under the simulated reading setting, \dn pregen60 reduces the mean full-ready latency to 5.299 seconds and the median latency to 0.027 seconds. This substantial reduction indicates that pre-generation changes the distribution of user-perceived latency rather than merely improving a small number of cases.

Compared with the image-continuation references, \dn pregen60 is faster on mean latency than StoryDiffusion (8.495 seconds) and IC-LoRA (30.057 seconds), while SDM-v2 remains faster at 2.827 seconds. This result should be interpreted with the output-scope distinction in mind: \dn returns a full playable response that includes both text and the corresponding scene image, whereas the image baselines report only image or workflow continuation latency. Thus, the result supports the claim that \dn's pre-generation mechanism can make full-ready interactive delivery competitive with specialized image-generation references under realistic player dwell time.

\paragraph{Effect of pre-generation.}
Table~\ref{tab:dn-pregen} isolates the mechanism-level effect of the simulated reading interval. Moving from normal20 to pregen60 reduces the mean full-ready latency from 27.095 seconds to 5.299 seconds and the median from 24.571 seconds to 0.027 seconds. The direct image rate remains 1.000 in both settings, and the placeholder rate is 0.0 in the source summary, indicating that the measured outputs are strict full-ready results rather than placeholder fallbacks. The likely-hit rate rises from 0.0 in normal20 to 0.8 in pregen60, which is consistent with the intended behavior of using reading time to prepare likely next-turn content.

\paragraph{Interpretation and limitations.}
The efficiency results provide a stable first piece of evidence for \dn's practical interaction profile. They show that the system is not merely generating richer multimodal content at arbitrary latency: when pre-generation has a realistic opportunity to operate, full-ready delivery becomes competitive with several image-only references. At the same time, the comparison is not a fully isomorphic system-to-system comparison. StoryDiffusion, SDM-v2, and IC-LoRA do not perform \dn's state update and textual continuation steps. We therefore use them as image-continuation efficiency references, while reserving separate sections for visual quality and narrative quality evaluation.

\paragraph{Stage-level conclusion.}
The current efficiency results are ready to support the first experimental subsection of the paper. They establish that \dn's pre-generation mechanism substantially improves next-turn delivery latency, reducing average full-ready latency by 80.4\% relative to the zero-dwell setting. The remaining quality question is orthogonal: the next sections will evaluate whether the content delivered quickly is also visually faithful, narratively coherent, and consistent across turns.

\section{Visual Quality Evaluation}
\label{sec:visual-quality}

\paragraph{Evaluation goal.}
This section focuses on whether the delivered images are usable for narrative play, not only on speed. To avoid fragmented reporting, visual findings are now represented through the Image Consistency rows in Table~\ref{tab:main-results-3in1}.

\paragraph{Current evidence and reporting policy.}
We keep two image-consistency channels in the unified table schema: an LLM-based score (Qwen3) and a feature-model score (DINO-v2 cosine). Both are currently marked as TODO because the cross-system runs still need to be frozen under a single evaluator pipeline. Until that freeze is complete, we avoid introducing additional standalone visual tables and keep all visual claims scoped to the same main-results structure.

\paragraph{Interpretation.}
This consolidation ensures that visual consistency is interpreted side by side with latency and text consistency, so later updates can be compared on a common baseline axis rather than across heterogeneous table layouts.

\section{Narrative and Text Quality Evaluation}
\label{sec:narrative-quality}

\paragraph{Evaluation goal.}
The text-side question is whether \dn keeps coherent and stable narrative progression across turns. We therefore align text-quality reporting with the same unified table used for efficiency and visual consistency.

\paragraph{Current evidence and reporting policy.}
Text consistency is represented in Table~\ref{tab:main-results-3in1} with two slots: Qwen3-based judged consistency (to be filled) and model-based consistency (currently N/A for these baselines because no shared domain-specific automatic scorer is frozen yet). This avoids mixing partially aligned subsets into separate headline tables and keeps the comparison contract explicit.

\paragraph{Interpretation.}
By reporting text consistency in the same matrix as latency and image consistency, the paper can discuss trade-offs across modalities directly and update remaining TODO entries without changing the experimental narrative structure.

\section{Ablation and Mechanism Analysis}
\label{sec:ablation}

\paragraph{Evaluation scope.}
The ablation study isolates whether key mechanisms contribute to both responsiveness and consistency. We compare the full system against variants that remove one mechanism at a time.

\begin{table*}[t]
\centering
\small
\resizebox{\textwidth}{!}{%
\begin{tabular}{lcccc}
\toprule
Variant & Mean Latency (s) & P95 Latency (s) & Image Consistency (Qwen3, 1--5) & Text Consistency (Qwen3, 1--5) \\
\midrule
Full model & 5.299 & 27.150 & TODO & TODO \\
$-$ pregen & 27.095 & 46.196 & TODO & TODO \\
$-$ council & TODO & TODO & TODO & TODO \\
$-$ readwait & TODO & TODO & TODO & TODO \\
\bottomrule
\end{tabular}
}
\caption{Ablation summary for mechanism contribution. The table isolates the effect of pregeneration, council planning, and read-wait strategy under aligned metrics. Pending artifact path: \texttt{[PATH\_TO\_ABLATION\_SUMMARY\_CSV]}.}
\label{tab:ablation-main}
\end{table*}

\paragraph{Interpretation.}
Table~\ref{tab:ablation-main} keeps the ablation axis directly aligned with the main-results metrics so each mechanism can be discussed in terms of both efficiency and consistency. The currently filled rows already show the dominant contribution of pregeneration to latency reduction, while remaining TODO entries are reserved for the same evaluation contract once the council and read-wait reruns are finalized.


\section{Conclusion}
The final conclusion will be completed after the quality-evaluation sections are integrated. The current efficiency results already show that pre-generation substantially reduces \dn's full-ready latency under a realistic reading-time simulation.

\bibliography{custom}

\appendix
\section{Experiment Artifact Provenance}
\label{app:provenance}
The efficiency tables in Section~\ref{sec:efficiency} use the latest v15 composite20 artifacts. Earlier exploratory runs, including v14, fixed-v2, and trimmed 16-sample variants, are retained only for internal analysis and are not used in the main paper numbers.

\begin{table*}[t]
\centering
\small
\resizebox{\textwidth}{!}{%
\begin{tabular}{lll}
\toprule
Artifact & Role in Paper & Path \\
\midrule
strict full-ready comparison & Main cross-system efficiency table & \texttt{experiments/handoff\_2026-04-30\_main\_experiments/01\_main\_tables/strict\_fullready\_main\_comparison\_2026-05-01\_v6\_composite20.csv} \\
DN normal vs. pregen60 & DN mechanism table & \texttt{experiments/handoff\_2026-04-30\_main\_experiments/01\_main\_tables/dn\_fullready\_normal\_vs\_pregen60\_2026-05-01\_v4\_composite20.csv} \\
v15 full-ready summary & Summary validation and mechanism signals & \texttt{experiments/handoff\_2026-04-30\_main\_experiments/01\_main\_tables/benchmark\_v15\_fullready\_nextturn\_20\_summary\_composite20\_v1.json} \\
per-sample v15 table & Per-item appendix candidate & \texttt{experiments/benchmark/standard\_runs/benchmark\_v15\_fullready\_nextturn\_20\_table\_composite20\_v1.csv} \\
\bottomrule
\end{tabular}
}
\caption{Source-of-truth artifacts for the efficiency evaluation. Paths are repository-relative.}
\label{tab:artifact-provenance}
\end{table*}

\section{Per-Sample Efficiency Details}
A future appendix revision will include per-sample latency rows from the v15 composite20 table to support sensitivity analysis and long-tail inspection.

\end{document}
